\itemhead{4.2 思考 · 未编号题组（教材p.110）}{XJ-B1-C04-TH-P110-Y1359}
\label{XJ-B1-C04-TH-P110-Y1359}
{\footnotesize 湘教版教材p.110 / PDF第117页 / 4.2 思考 / 未编号题组}
\field{原题}{当 $a>1$ 时，$y=a^x$ 和 $y=\left(\dfrac1a\right)^x$ 的图象有什么关系？能否利用 $y=a^x$ 的图象画出 $y=\left(\dfrac1a\right)^x$ 的图象？}
\field{参考答案}{两图象关于 $y$ 轴对称，因为 $(1/a)^x=a^{-x}$；将每个点 $(x,a^x)$ 映到 $(-x,a^x)$ 即得。}
{\small \field{关联知识}{E01 指数函数的概念、图象与性质（教材帮p.150、151 / PDF159、160）；F09 图象平移、伸缩与对称（教材帮p.86、94 / PDF95、103）}
\field{方法与易错点}{先查底数条件，利用已知点定底数，再据单调性及图象求值或范围。 易错：只要求底数非负；漏排1；将 $a^x+b$ 当标准指数函数且仍用原值域。}
\field{答案依据}{答案册未提供本栏目的独立对应答案；依据原题、教材概念补写参考回应。开放性总结不设唯一数值答案。}}
